%% Um 90° gedrehte Tabelle
%
\begin{sidewaystable}[p]
\raggedleft%% Rechtsbündig, damit die Tabelle oben ausgerichtet wird (Forderung KSP)
\renewcommand{\arraystretch}{1.3}%% Etwas mehr Abstand zwischen den Tabellenzeilen
%% Tabelle mit einer festen Breite (hier: so breit wie die Seitenhöhe).
%% Spaltendefinition:
%% Erste Spalte vom Typ "l": Spalte so breit, wie der Inhalt, linksbündige Ausrichtung
%% Zweite Spalte vom Typ "L": Spalte mit einer festen Breitenangabe, linksbündiger Flattersatz. Inhalt wird umbrochen, wenn zu lang
%% 2 Spalten vom Typ "x":  die Breite wird automatisch berechnet
%% in Abhängigkeit von der Gesamtbreite der Tabelle und Breite anderer Spalten,
%% Inhalt wird umbrochen, wenn zu lang
%% Format: linksbündiger Flattersatz ohne Silbentrennung, 
%% Weitere Möglichkeiten:
%% "X" = Blocksatz (mit Silbentrennung)
%% "F" = linksbündiger Flattersatz mit Silbentrennung (KSP-Verlag: unerwünscht),
%% "y" = rechtsbündiger Flattersatz ohne Silbertrennung
%% "Y" = rechtsbündiger Flattersatz mit Silbertrennung
\begin{tabularx}{\linewidth}{l L{10em} *{2}x}%
	\toprule%
	\rowcolor{tableheadcolor}%
	\tableheadfont Spalte vom Typ \enquote{l} &	
	\tableheadfont Spalte vom Typ \enquote{L} & 
	\tableheadfont Spalte vom Typ \enquote{x} & 
	\tableheadfont Spalte vom Typ \enquote{x}
	\tabularnewline%
	\midrule%
		Spaltenbreite = Breite des Inhalts & 
		Spalte mit der festen Breite von 10\,em & 
		Spaltenbreite wird automatisch bestimmt. Inhalt wird bei Bedarf umbrochen. & 
		Spaltenbreite wird automatisch bestimmt. Inhalt wird bei Bedarf umbrochen.
	\tabularnewline
	%%-----------------------------------------------------------
	Nächste Zeile & Etwas längerer Inhalt & Weiterer Inhalt & und noch etwas \tabularnewline
	Weitere Zeile & Etwas längerer Inhalt & Weiterer Inhalt & und noch etwas \tabularnewline
	%% Zwischenkopf ---------------------------------------------
	\rowcolor{tablesubheadcolor}%
	\multicolumn{4}{l}{\bfseries\color{KITblue} Zwischenkopfzeile, welche über mehrere Spalten geht}%
	\tabularnewline%
	%%-----------------------------------------------------------
	Inhalt  & Etwas längerer Inhalt & Weiterer Inhalt & und noch etwas \tabularnewline
	Inhalt  & Etwas längerer Inhalt & Weiterer Inhalt & und noch etwas \tabularnewline
	Inhalt  & Etwas längerer Inhalt & Weiterer Inhalt & und noch etwas \tabularnewline
	%% Zwischenkopf ---------------------------------------------
	\rowcolor{tablesubheadcolor}%
	\multicolumn{4}{l}{\bfseries\color{KITgreen} Zwischenkopfzeile, welche über mehrere Spalten geht}%
	\tabularnewline
	%%-----------------------------------------------------------
	Inhalt  & Etwas längerer Inhalt & Weiterer Inhalt & und noch etwas \tabularnewline
	Inhalt  & Etwas längerer Inhalt & Weiterer Inhalt & und noch etwas \tabularnewline
	\bottomrule%
\end{tabularx}%
\caption{Beispiel für eine breite, gedrehte Tabelle}%
\label{tab:ex-sideways}%
\end{sidewaystable}